\section{Positive integral ranges}\label{sec:positive}

The examples concern small degrees for fixed embeddings. Sufficiently high degree gives an integral theorem. The following explicit regularity bound is a convenient form of the sufficiently ample argument of \citet[Lemma 2.12, pp.~242--243]{BMS}; it is not a claim that the high-degree theorem is new. Regularity below is Castelnuovo--Mumford regularity for the specified embedding, with coherent sheaves on the variety viewed as sheaves on its ambient projective space.

\begin{theorem}\label{thm:regularity}
Let $V\subset\PP^r_k$ be smooth projective of dimension $n\ge2$, and put
\[
 d_{\mathrm{int}}(V)=\max\bigl(1,\operatorname{reg}(\Omega_V^a):0\le a\le n\bigr).
\]
If $d\ge d_{\mathrm{int}}(V)$ and $Y\in|\OO_V(d)|$ is smooth, then crystalline restriction is an isomorphism below degree $n-1$ and injective with torsion-free cokernel in degree $n-1$. In that degree it identifies the full torsion subgroups. Coherent restriction in degree $n-1$ is injective already for
$d\ge\max(1,\operatorname{reg}\omega_V)$.
\end{theorem}
\begin{proof}
For every $t\ge d_{\mathrm{int}}(V)$, regularity gives
$H^b(V,\Omega_V^a(t))=0$ for $b>0$. By Serre duality and the wedge pairing,
\[
 H^b(V,\Omega_V^a(-t))=0\qquad(b<n).
\]
In particular this holds for $t=md$, $m\ge1$. Write $D^{\ge q}$ for complexes with zero cohomology in degrees below $q$. The divisor sequence then shows
\[
 R\Gamma(Y,\Omega_V^a|_Y(-md))\in D^{\ge n-1}\quad(m\ge1),
\]
and the cone of $R\Gamma(\Omega_V^a)\to R\Gamma(\Omega_V^a|_Y)$ is in $D^{\ge n-1}$.

For $0\le a\le n-1$, exterior powers of the conormal sequence give the exact resolution
\[
 0\to\OO_Y(-ad)\to\Omega_V^1|_Y(-(a-1)d)\to\cdots
 \to\Omega_V^{a-1}|_Y(-d)\to\Omega_V^a|_Y\to\Omega_Y^a\to0.
\]
Each of its $a$ preceding terms has cohomology in $D^{\ge n-1}$. Successive kernels therefore show
\[
 \operatorname{Cone}\bigl(R\Gamma(V,\Omega_V^a)
                   \to R\Gamma(Y,\Omega_Y^a)\bigr)
 \in D^{\ge n-a-1}.
\]
For $a=n$, the same assertion is immediate from $\Omega_Y^n=0$. Shifting each form-degree piece by $[-a]$ and using the finite de Rham filtration puts the de Rham restriction cone in $D^{\ge n-1}$.

Let $K$ be the crystalline restriction cone. It is a bounded complex with finite cohomology over $W$, and derived crystalline base change identifies $K\otimes_W^{\mathbf L}k$ with the de Rham cone. Its universal-coefficient sequences are
\[
 0\to H^j(K)/p\to H^j(K\otimes_W^{\mathbf L}k)
                 \to H^{j+1}(K)[p]\to0.
\]
Nakayama gives $H^j(K)=0$ for $j<n-1$. Taking $j=n-2$ also gives $H^{n-1}(K)[p]=0$, so that group is torsion-free. The long exact restriction sequence gives the lower-degree isomorphisms and the degree-$(n-1)$ injection. Its cokernel embeds into $H^{n-1}(K)$ and is therefore torsion-free. An injection of finite $W$-modules with torsion-free quotient identifies their torsion subgroups.

Finally, Serre duality gives
$H^{n-1}(V,\OO_V(-d))^\vee=H^1(V,\omega_V(d))$. The stated regularity bound kills this group, and the divisor sequence gives coherent injectivity.
\end{proof}

There are also familiar hypotheses giving the full range $d\ge1$. For a smooth global complete intersection $V\subset\PP^r$, every smooth hypersurface section has the integral conclusion of Theorem~\ref{thm:regularity}. The same holds if $n\le p$ and $V$ lifts smoothly to $W_2(k)$. Without that dimension bound, a $W_2$-lift gives isomorphism below degree $\min(n,p)-1$ and injection with torsion-free cokernel in that degree. These follow from the Hodge-equivalence statements of \citet[Definition 2.2.1 and Theorem 2.2.3, pp.~1854--1855]{CotnerZavyalov} by the de Rham filtration and universal-coefficient argument just proved. For complete intersections apply their comparison with projective space to $V$ and to $Y$ and use the triangle of the two cones. The lifting statement does not require a lift of the chosen divisor. None of these hypotheses is supplied by smoothness and very ampleness alone.
